\section{Introduction}

In the frictionless Black--Scholes--Merton model, continuous trading in the
underlying and money-market account permits exact replication of a European
option under the standard assumptions \citep{blackscholes1973,merton1973}.
Implemented strategies instead trade on a finite grid and may incur transaction
costs. The resulting hedging problem balances residual risk against trading
activity and is generally different from continuous-time frictionless
replication. Analytical and numerical approaches remain available in structured
frictional models, including asymptotic approximations, stochastic control, and
variational inequalities
\citep{leland1985,hoggardwhalleywilmott1994,whalleywilmott1997}.

\citet{buehler2019} formulate hedging with generic frictions as a discrete-time
stochastic-control problem and approximate admissible strategies with neural
networks. Their framework supports nonlinear risk criteria and multiple hedging
instruments. The present work does not attempt to reproduce their complete
architecture or empirical design. It studies a smaller setting in which a
shared-weight NumPy network hedges a European call under empirical CVaR.

The aims of this paper are limited and technical:

\begin{enumerate}[leftmargin=2em]
  \item state the loss, trading-cost convention, empirical-CVaR estimator, and
  analytic derivatives in a form that can be checked line by line;
  \item compare the manual derivatives with centered finite differences and
  document which validation steps are complete or still outstanding;
  \item preserve the archived numerical results while describing them as
  single-run observations rather than evidence of statistical dominance; and
  \item separate the present state-only, stock-only experiments from the
  broader Deep Hedging framework on which they are based.
\end{enumerate}

The paper is best read as an implementation and reproducibility study. It does
not propose a new deep-hedging algorithm, establish optimality, or claim that a
neural policy generally dominates classical hedging rules.

\section{Scope Relative to the Deep Hedging Literature}
\label{sec:scope}

Table~\ref{tab:scope} summarizes the main differences between the reported
experiment in \citet{buehler2019} and the present implementation. The comparison
is included to prevent a qualitative resemblance from being mistaken for a
full numerical reproduction.

\begin{table}[htbp]
\centering
\caption{Scope of the present implementation relative to the original Deep
Hedging experiment.}
\label{tab:scope}
\begin{tabular}{>{\raggedright\arraybackslash}p{0.22\textwidth}
                >{\raggedright\arraybackslash}p{0.31\textwidth}
                >{\raggedright\arraybackslash}p{0.34\textwidth}}
\toprule
Component & \citet{buehler2019} & Present study \\
\midrule
Policy state
& Market information and the preceding position enter the semi-recurrent policy.
& Time and log-moneyness enter the GBM policy; truncated variance is added under
Heston. The previous position is absent from the reported state-only results. \\
\addlinespace
Time parametrization
& The reported numerical study uses date-specific parameters; recurrent variants
are also discussed.
& One shared-weight network is evaluated at every trading date. \\
\addlinespace
Network
& The reported implementation uses two hidden layers, ReLU activations, and batch
normalization.
& Two hidden layers with 16 $\tanh$ units each and no batch normalization. \\
\addlinespace
Heston instruments
& Stock and an idealized variance swap.
& Stock only. \\
\addlinespace
Heston comparison
& A model-consistent complete-market comparison using both instruments.
& A fixed-volatility Black--Scholes delta with less information than the neural
policy. \\
\addlinespace
Primary criterion
& General convex risk measures and indifference-pricing experiments.
& Empirical $\CVaR_{0.95}$ of terminal loss under a fixed cash convention. \\
\bottomrule
\end{tabular}
\end{table}

Accordingly, the numerical section examines whether the fitted policy displays
lower trading intensity and different tail-loss point estimates in the chosen
simulations. It does not test equivalence to the original architecture,
instrument set, or Heston comparison.
